% This is samplepaper.tex, a sample chapter demonstrating the
% LLNCS macro package for Springer Computer Science proceedings;
% Version 2.21 of 2022/01/12
%
\documentclass[runningheads]{llncs}
%
\usepackage[T1]{fontenc}
% T1 fonts will be used to generate the final print and online PDFs,
% so please use T1 fonts in your manuscript whenever possible.
% Other font encondings may result in incorrect characters.
%
\usepackage[utf8]{inputenc}
\usepackage{url}
\usepackage{graphicx}
\usepackage{hyperref}
% Used for displaying a sample figure. If possible, figure files should
% be included in EPS format.
%
% If you use the hyperref package, please uncomment the following two lines
% to display URLs in blue roman font according to Springer's eBook style:
%\usepackage{color}
%\renewcommand\UrlFont{\color{blue}\rmfamily}
%\urlstyle{rm}
%

\begin{document}
%
\title{Traqula: Providing a Foundation for The Evolving SPARQL Ecosystem Through Modular Query Parsing, Transformation, and Generation}
%
%\titlerunning{Abbreviated paper title}
% If the paper title is too long for the running head, you can set
% an abbreviated paper title here
%
\author{Jitse De Smet\inst{1}\orcidID{0009-0002-6513-5013} \and
Ruben Taelman\inst{1}\orcidID{0000-0001-5118-256X}}
%
\authorrunning{J. De Smet and R. Taelman}
% First names are abbreviated in the running head.
% If there are more than two authors, 'et al.' is used.
%
\institute{IDLab, Department of Electronics and Information Systems, Ghent University – imec.
\email{\{firstname\}.\{lastname\}@ugent.be}}
%
\maketitle              % typeset the header of the contribution
%
\begin{abstract}

    SPARQL engines continue to diverge in the versions and language extensions they support. With SPARQL 1.2 nearing completion and the working group’s plans toward future maintenance, this divergence is likely to grow. As a result, language tools that aim to support multiple SPARQL versions and dialects face increasing maintenance costs. Furthermore, areas such as SPARQL federation become harder to manage. To address these challenges, we introduce Traqula, a modular TypeScript-based framework for parsing, transforming, and generating SPARQL queries. This article presents Traqula’s architecture and its builder-based dependency-injection design, and we show how these components provide modular support for SPARQL 1.1 and 1.2 parsing and generation. Traqula enables round-tripping, supports generic algebraic and AST transformations, and delivers significantly faster parsing compared to current JavaScript-based SPARQL parsing. Traqula’s modular architecture enables flexible experimentation for researchers, and it is production-ready thanks to extensive testing and modern development techniques. Looking forward, Traqula lays the groundwork for supporting additional query languages, developing generic rewriting modules for SPARQL federation engines, and building robust tooling for an increasingly diverse SPARQL ecosystem.

\keywords{SPARQL, query \and algebra \and modularity \and parser \and generator}
\end{abstract}
%
%
%
\section{Introduction}
\label{introduction}


The SPARQL query language~\cite{spec:sparqllang} is the standard way to query RDF~\cite{spec:rdf} data.
While SPARQL endpoints~\cite{spec:sparqlprot} are commonly used to expose Knowledge Graphs through a highly expressive query API,
alternative APIs have been proposed~\cite{tpf,smartkg,sage,wisekg,brtpf,passage} that offer different trade-offs between client and server effort for query execution.
This heterogeneity in server APIs introduces challenges when executing federated SPARQL queries~\cite{spec:sparqlfed,fedx,hibiscus,splendid,hanski_wikidata_federation_2025} over multiple of these APIs.
While this API-based heterogeneity has been an active field of research~\cite{comunica,heterogeneous_lars,heterogeneous_fedqpl,heterogeneous_replicas} in recent years,
the heterogeneity of the SPARQL language itself lacks understanding.
In practice, many SPARQL dialects exist~\cite{geosparql,tsparql,csparql,rdfversioningquerylang} each introducing their own extensions or limitations.
Virtuoso, for example, extends SPARQL with full-text search capabilities \cite{virtuoso-full-text-search},
Apache Jena adds support for constructing quads \cite{jena-construct-quad},
and Oxigraph provides an additional built-in function, \texttt{ADJUST} \cite{sparql-adjust,oxigraph-adjust}.
Furthermore, some SPARQL endpoints might limit the SPARQL language
by removing expensive operators such as OPTIONAL \cite{querycomplexity}.
In the near future, this heterogeneity is expected to increase even further as the RDF and SPARQL W3C Working Group enters its
maintenance mode \cite{rdf12-charter,modular-parsing},
which could deliver more rapid successions of SPARQL once SPARQL 1.2 is finalized.
As a result, the growing diversity of SPARQL versions and dialects introduces several challenges:

\begin{enumerate}
    \item \textbf{Query evaluation}: A user query written in one version might not be executable by a SPARQL engine supporting another version.
    \item \textbf{Tooling}: Linters, formatters, editors, and language servers often assume a specific SPARQL version.
    \item \textbf{Maintainability}: Tools that support multiple SPARQL versions typically do so by maintaining multiple software versions, or one version with many conditions, making maintainability highly challenging.
\end{enumerate}

\begin{figure}[htbp]
    \centering
    \includegraphics[width=0.85\textwidth]{img/traqula-federation}
    \caption{The federated SPARQL query (blue) uses SPARQL 1.2 features such as triple terms and uses the non-standard builtin function ADJUST~\cite{sparql-adjust}.
    The query targets four SPARQL endpoints, all supporting different SPARQL version, RDF profiles~\cite{rdf-1-2}, and language extensions.
    To allow query engines to integrate this heterogeneity,
        frameworks such as Traqula are necessary to bridge between these SPARQL dialects.}
    \label{federation}
\end{figure}


The query evaluation problem within a heterogeneous SPARQL ecosystem becomes especially visible in federated SPARQL query execution.
While language dialects are common in many technologies, SQL~\cite{iso-sql} being a well-known example,
the RDF data model~\cite{spec:rdf} is explicitly designed for seamless integration of distributed datasets.
SPARQL reflects this distributed nature through support for federated queries,
where a single query may involve multiple endpoints, each with its own capabilities, limitations, and language features.

In such a setting, a SPARQL query written in one SPARQL version or dialect may not be executable on all federation members.
For instance, a query formulated in SPARQL 1.2 might rely on features, functions,
or syntactic constructs that an older endpoint does not support, or that are only available as vendor-specific extensions.
\autoref{federation} illustrates this scenario: a query is written in SPARQL 1.2 and uses the non-standard ADJUST function~\cite{sparql-adjust},
yet the endpoints differ in their supported SPARQL versions,
RDF profiles~\cite{rdf-1-2}, and feature sets.
Although the SPARQL Service Description specification~\cite{spec:sparqlservicedesc} allows endpoints to declare their supported features,
it does not offer mitigation paths to resolve these language mismatches.
To solve this problem, there is a need for a parsing, transformation, and generation framework such as Traqula that can handle various SPARQL dialects.
Such a framework provides a foundation for future SPARQL federation research towards new techniques and algorithms to manage these dialects.

In prior work~\cite{modular-parsing},
we introduced our vision of a modular SPARQL parser to address the growing heterogeneity of SPARQL dialects using
builder-based dependency injection~\cite{fowler-dep-inj}.
We envisioned a design with a prototype implementation demonstrating parser composability,
showing how grammar modules could be combined to modularly define multiple SPARQL parsers covering different versions and dialects.
In this article, we fully realize that vision with a complete implementation
as well as applying the modular architecture to query generation and transformation.
These features allow queries to be parsed, transformed, and regenerated reliably, while preserving their structure and semantics across different dialects.

This article presents Traqula, a modular SPARQL toolkit implementing these ideas.
Traqula has already been integrated into the widely used Comunica SPARQL querying framework~\cite{comunica}, demonstrating its maturity and practical applicability.
Traqula's modular and composable architecture enables researchers and practitioners to:
\begin{enumerate}
    \item Experiment with grammar changes, including adding, modifying, or removing rules;
    \item tackle the complexities of federated SPARQL querying in a heterogeneous SPARQL ecosystem; and
    \item lay the foundation for future query formatting and rewriting tools.
\end{enumerate}
Traqula is implemented in TypeScript and is available as open-source software under the MIT license on
GitHub and npm,
providing the community with a robust and flexible resource for building tools and engines for a heterogeneous SPARQL environment.


The remainder of this paper is structured as follows.
\autoref{requirements} details the requirements of Traqula, \autoref{sota} reviews related work,
\autoref{architecture} presents the high-level architecture of Traqula,
and \autoref{implementation} discusses its implementation.
\autoref{perf} evaluates its performance, and \autoref{conclusion} concludes the paper.


%
% the environments 'definition', 'lemma', 'proposition', 'corollary',
% 'remark', and 'example' are defined in the LLNCS documentclass as well.
%
\begin{proof}
Proofs, examples, and remarks have the initial word in italics,
while the following text appears in normal font.
\end{proof}
For citations of references, we prefer the use of square brackets
and consecutive numbers. Citations using labels or the author/year
convention are also acceptable. The following bibliography provides
a sample reference list with entries for journal
articles~\cite{ref_article1}, an LNCS chapter~\cite{ref_lncs1}, a
book~\cite{ref_book1}, proceedings without editors~\cite{ref_proc1},
and a homepage~\cite{ref_url1}. Multiple citations are grouped
\cite{ref_article1,ref_lncs1,ref_book1},
\cite{ref_article1,ref_book1,ref_proc1,ref_url1}.

\begin{credits}
\subsubsection{\ackname} A bold run-in heading in small font size at the end of the paper is
used for general acknowledgments, for example: This study was funded
by X (grant number Y).

\subsubsection{\discintname}
It is now necessary to declare any competing interests or to specifically
state that the authors have no competing interests. Please place the
statement with a bold run-in heading in small font size beneath the
(optional) acknowledgments\footnote{If EquinOCS, our proceedings submission
system, is used, then the disclaimer can be provided directly in the system.},
for example: The authors have no competing interests to declare that are
relevant to the content of this article. Or: Author A has received research
grants from Company W. Author B has received a speaker honorarium from
Company X and owns stock in Company Y. Author C is a member of committee Z.
\end{credits}
%
% ---- Bibliography ----
%
% BibTeX users should specify bibliography style 'splncs04'.
% References will then be sorted and formatted in the correct style.
%
\bibliographystyle{splncs04}
\bibliography{references}
\end{document}
